% 三步搜索教学构造：搜索范围p=7，步长4、2、1，代价最小处依次为(4,4)、(6,6)、(7,6)；图像坐标纵轴向下。宽约95mm。
\begin{tikzpicture}[x=4mm,y=-4mm,font=\figurefont\small,
  lab/.style={font=\figurefont\footnotesize,text=FigureMuted},
  path/.style={-{Stealth[length=1.8mm]},draw=FigureStroke,line width=.8pt}]
  \draw[FigureGrid,line width=.4pt] (-7.5,-7.5) rectangle (7.5,7.5);
  \foreach \x in {-7,...,7}{
    \foreach \y in {-7,...,7}{
      \fill[FigureGrid] (\x,\y) circle (.6pt);
    }
  }
  \foreach \x in {-4,0,4}{
    \foreach \y in {-4,0,4}{
      \fill[FigureBlue] (\x,\y) circle (2.4pt);
    }
  }
  \foreach \x in {2,4,6}{
    \foreach \y in {2,4,6}{
      \draw[FigureYellow,fill=FigureYellowFill,line width=.9pt] (\x,\y) circle (2.6pt);
    }
  }
  \foreach \x in {5,6,7}{
    \foreach \y in {5,6,7}{
      \draw[FigureRed,line width=.9pt] ({\x-.22},{\y-.22}) rectangle ({\x+.22},{\y+.22});
    }
  }
  \draw[path] (0,0) -- (3.75,3.75);
  \draw[path] (4,4) -- (5.8,5.8);
  \draw[path] (6,6) -- (6.75,6);
  \fill[FigureRed] (7,6) circle (2pt);
  \draw[-{Stealth[length=1.8mm]},FigureStroke,line width=.6pt] (-8.5,-8.3) -- (-6.5,-8.3) node[right,lab] {$d_x$};
  \draw[-{Stealth[length=1.8mm]},FigureStroke,line width=.6pt] (-8.5,-8.3) -- (-8.5,-6.3) node[below,lab] {$d_y$};
  \fill[FigureBlue] (9.5,-5) circle (2.4pt);
  \node[anchor=west] at (10,-5) {第1步，步长4};
  \draw[FigureYellow,fill=FigureYellowFill,line width=.9pt] (9.5,-3) circle (2.6pt);
  \node[anchor=west] at (10,-3) {第2步，步长2};
  \draw[FigureRed,line width=.9pt] (9.28,-1.22) rectangle (9.72,-.78);
  \node[anchor=west] at (10,-1) {第3步，步长1};
  \fill[FigureRed] (9.5,1) circle (2pt);
  \node[anchor=west] at (10,1) {结果$(7,6)$};
  \node[anchor=west,lab,align=left] at (9.3,4.2) {灰点：全搜索的\\$15\times15=225$个候选};
\end{tikzpicture}
